\documentclass[11pt,a4paper]{article}
\usepackage[T2A]{fontenc}
\usepackage[utf8]{inputenc}
\usepackage[russian]{babel}
\usepackage[margin=2.2cm]{geometry}
\usepackage{amsmath,amssymb}
\usepackage{booktabs}
\usepackage{array}
\usepackage{enumitem}
% xcolor нужен до hyperref: тот красит ссылки выражением blue!50!black,
% а смешение цветов через «!» умеет только xcolor (не базовый color).
\usepackage{xcolor}
\usepackage[colorlinks=true,linkcolor=blue!50!black,urlcolor=blue!50!black]{hyperref}
\usepackage{parskip}

\newcommand{\med}{\operatorname{med}}

\title{Темп прохождения курса:\\ что означают числа в таблице участников}
\author{standings-edu}
\date{16 сентября 2026}

\begin{document}
\maketitle

\section{О чём этот документ}

На странице \textbf{«Статистика участников»} (и в профиле ученика) каждому
ученику показывается строка чисел: прогресс, скорость, часы, прогноз. Здесь
выписано, откуда берётся каждое из них --- формулой и словами.

Документ описывает \emph{реализованную} модель: все формулы ниже соответствуют
коду в \texttt{course\_stats.go} (полные пути --- в последнем разделе). Числа
считаются при генерации из уже собранных посылок; отдельных запросов к сайтам
не делается.

\textbf{Курс} --- это упорядоченный список задач группы: контесты страницы
снизу вверх, внутри контеста --- слева направо. Обозначения: $i$ --- ученик,
$j$ --- задача курса, $M$ --- число задач, $S_i$ --- решённые учеником задачи
курса.

\subsection*{Сводка: колонка $\to$ формула}

\begin{center}
\begin{tabular}{@{}lp{5.4cm}p{6.1cm}@{}}
\toprule
Колонка & Что означает & Формула \\
\midrule
Прогресс & доля пройденного по сложности &
  $P_i = \dfrac{\sum_{j \in S_i} w_j}{\sum_{j} w_j}$ \\[2.2ex]
Скорость & темп относительно группы &
  $v_i = \dfrac{\sum w_j}{\sum \max(T_{ij},\, w_j/2)}$,
  делённое на медиану группы \\[2.2ex]
По времени & та же скорость, свежее весит больше &
  то же с весом сессии $\gamma_s = 2^{-d_s/28}$ \\[1.2ex]
Часов & активное время на задачах курса &
  $A_i = \sum_{j} T_{ij}$ \\[1.2ex]
$\approx$ч/нед & типичная недельная занятость &
  медиана по неделям с активностью (из 8) \\[1.2ex]
$\approx$недель & прогноз до конца курса &
  $W_{\text{ост}} \,/\, (v^{\text{тек}}_i \cdot E_i \cdot h_i)$ \\[1.2ex]
Сейчас на & фронт & последняя решённая по порядку курса \\[1.2ex]
Сигналы & застрял / пропустил &
  $T_{ij}/w_j > 3$ \ \ / \ \ дальше решено $\geq 2$ \\[1.2ex]
Флаги & подозрительные паттерны посылок & четыре детектора, разд.~\ref{sec:flags} \\
\bottomrule
\end{tabular}
\end{center}

Всё держится на двух величинах: \textbf{активное время} $T_{ij}$ (сколько
ученик реально сидел над задачей) и \textbf{вес} $w_j$ (сколько над ней сидит
типичный ученик). С них и начнём.

\section{Активное время: отделяем работу от перерывов}

Посылки ученика сортируем по времени: $t_1 < t_2 < \dots$ Если пауза между
соседними посылками больше порога, начинается новая \textbf{сессия} (занятие):
\[
  t_k - t_{k-1} > \tau, \qquad \tau = 45\ \text{мин}.
\]
Паузы между посылками бимодальны --- минуты внутри занятия и часы-дни между
занятиями, --- поэтому результат не слишком чувствителен к выбору $\tau$.
Перерывы любой длины (выходные, каникулы, болезнь) в активное время
\textbf{не попадают вовсе}: «не решал месяц» $\neq$ «решал задачу месяц».

Внутри сессии пауза перед посылкой приписывается задаче \emph{этой} посылки ---
это время, потраченное на её подготовку. Первая посылка сессии получает
надбавку $\delta_0$ на ненаблюдаемый «вход» (чтение условия, написание кода):
\[
  \delta_0 = \min\bigl(\med(\text{внутрисессионных пауз ученика}),\ 10\ \text{мин}\bigr).
\]

Активное время ученика $i$ на задаче $j$:
\[
  T_{ij} \;=\; \sum_{k:\ \text{посылка }k\text{ по задаче }j}
      \underbrace{\bigl(t_k - t_{k-1}\bigr)}_{\text{если }k\text{ не первая в сессии}}
  \;+\; \delta_0 \cdot [\,\text{$j$ открывает сессию}\,].
\]

\textbf{Читается так:} время между двумя соседними посылками записывается на
задачу второй из них; первая посылка занятия получает $\delta_0$.

Сессии строятся по \emph{всем} посылкам ученика (сессия --- свойство человека,
а не курса), но в колонку \textbf{«Часов»} идёт только время задач этого курса:
\[
  A_i \;=\; \sum_{j = 1}^{M} T_{ij}.
\]
Если ученик параллельно занимается в другой группе, её время в $A_i$ не входит.

\section{Вес задачи: сколько она «стоит»}

Вес $w_j$ --- типичное время решения задачи в минутах. Оценивается из данных
когорты в три шага.

\textbf{Шаг 1. Сырая оценка} --- медиана активного времени тех, кто решил:
\[
  \widehat{w}_j = \med\bigl\{\, T_{ij} \;:\; \text{ученик } i \text{ решил } j,\ T_{ij} > 0 \,\bigr\},
  \qquad n_j = \bigl|\{\,i\,\}\bigr|.
\]
Медиана, а не среднее --- чтобы один застрявший не раздул вес.

\textbf{Шаг 2. Сглаживание к типичной задаче курса.} Пока решивших мало,
оценка шумная, поэтому её притягивают к средней задаче:
\[
  w_j \;=\; \frac{n_j\,\widehat{w}_j + n_0\,\bar{w}}{n_j + n_0},
  \qquad n_0 = 5,
\]
где $\bar{w}$ --- медиана \emph{ненулевых} $\widehat{w}_j$ по курсу (если
данных нет вовсе --- условные 20 минут). При $n_j = 0$ получается
$w_j = \bar w$, то есть «обычная задача»; с накоплением данных вес сам
уточняется при каждой генерации.

\textbf{Шаг 3. Поправка на редкость решения.} Медиана времени \emph{решивших}
занижает сложность редких задач: их решают немногие --- сильные и быстрые.
Поэтому вес умножается на редкость среди \emph{дошедших}:
\[
  w_j \;\mathrel{{*}{=}}\; \min\!\Bigl(c_{\max},\ \bigl(\tfrac{1}{p_j}\bigr)^{\beta}\Bigr),
  \qquad p_j = \frac{\text{решивших } j}{\text{дошедших до } j},
  \qquad \beta = \tfrac12,\ c_{\max} = 3.
\]
«Дошёл» --- фронт ученика (последняя решённая по порядку курса) не раньше
задачи $j$; ученик без единой решённой не дошёл никуда. Нормировка по дошедшим,
а не по всей группе --- иначе поздние задачи дорожали бы просто за позицию в
курсе. При числе дошедших меньше 5 доля слишком шумная, и поправки нет.

\textbf{Читается так:} если до задачи дошли 20 человек, а решили 5, то
$p = 0{,}25$ и вес умножается на $\sqrt{4} = 2$.

\subsection*{Чем это отличается от «типичного времени» в таблице контеста}

В таблице контеста у задачи (по доступу с просмотром) показывается
\textbf{сырая медиана} $\widehat{w}_j$ --- без сглаживания и без поправки, и
только если решивших с известным временем хотя бы 3. Это честное «типичное
время решивших». В расчётах темпа участвует полный вес $w_j$ после всех трёх
шагов.

\section{Прогресс}

\[
  P_i \;=\; \frac{\sum_{j \in S_i} w_j}{\sum_{j=1}^{M} w_j} \;\in\; [0,1].
\]

\textbf{Читается так:} доля пройденного \emph{по сложности}, а не по числу
задач: две простые задачи весят меньше одной трудной. Рядом в интерфейсе ---
«решено~/~всего» в штуках. Задачи, решённые на сайтах без времени посылок
(ACMP), в прогресс входят как обычные.

\section{Скорость}
\label{sec:speed}

Пусть $S_i^{+}$ --- решённые задачи курса, у которых есть активное время
($T_{ij} > 0$). Сырая скорость:
\[
  v_i \;=\; \frac{\sum_{j \in S_i^{+}} w_j}
                 {\sum_{j \in S_i^{+}} \max\bigl(T_{ij},\ \alpha\, w_j\bigr)},
  \qquad \alpha = \tfrac12 .
\]

\textbf{Читается так:} сколько минут \emph{типичного} времени ученик закрывает
за минуту \emph{своего}. Числитель --- объём материала в единицах типичного
времени, знаменатель --- личное время на тех же задачах.

Три решения, которые стоит понимать:

\begin{itemize}[itemsep=2pt]
  \item \textbf{Только решённые задачи в знаменателе.} Время, утопленное в
        нерешённых, скорость не занижает --- оно видно отдельно, в активных
        часах и в сигнале «застрял». Иначе ученик, честно бьющийся над
        сложным, выглядел бы медленнее того, кто сложное просто пропускает.
  \item \textbf{Пол $\alpha w_j$.} Решённая задача не может «стоить» меньше
        половины своего типичного времени. Обдумывание до первой посылки
        модель не видит (кредитуется только $\delta_0$), и без пола решение
        одной посылкой давало бы фиктивно высокий темп. Пол ограничивает вклад
        одной задачи величиной $\times(1/\alpha) = \times 2$ от типичного.
  \item \textbf{Задачи без времени не участвуют вовсе.} Сайт без времени
        посылок (ACMP) или эпизод, исключённый флагом честности, дал бы вес в
        числитель, не дав ни минуты в знаменатель, --- и раздул бы $v_i$.
        Такие задачи выкидываются из обеих частей (но остаются в прогрессе).
\end{itemize}

\subsection*{Нормировка на группу}

Личное время на задаче обычно правее медианы (распределение скошено), поэтому
сырая $v_i$ у большинства меньше единицы. Показываемое число
перецентрируется на когорту:
\[
  v_i^{\text{норм}} \;=\; \frac{v_i}{\med_k v_k},
\]
по ученикам с достаточными данными, и только если таких хотя бы 5 (иначе
остаётся сырая шкала). Медианный ученик получает ровно $\times 1$, ранжирование
не меняется. Тем же множителем нормируется и «текущая форма».

\textbf{Читается так:} $\times 1$ --- как типичный ученик когорты,
$\times 2$ --- вдвое быстрее, $\times 0{,}5$ --- вдвое медленнее.

\subsection*{Когда число не показывается}

При $A_i < 2$~ч или $|S_i| < 5$ вместо скорости выводится «мало данных».

\subsection*{Относительно какой когорты}

В таблице группы веса задач и медиана для нормировки считаются \emph{по составу
этой группы}. В профиле ученика есть переключатель «все группы»: там то же
самое, но по объединению составов всех групп, где есть эти контесты. Поэтому
$\times 1$ в двух видах --- это две разные базы сравнения.

\section{По времени: текущая форма}

Та же скорость, но свежие занятия весят больше. Сессия, закончившаяся $d_s$
дней назад, получает вес с полупериодом $H = 28$ дней:
\[
  \gamma_s = 2^{-d_s/H},
  \qquad
  v_i^{\text{тек}} \;=\;
  \frac{\sum_s \gamma_s \sum_{j \text{ решена в } s} w_j}
       {\sum_s \gamma_s \sum_{j \in S_i^{+}} q_{sj}\cdot\frac{\max(T_{ij},\,\alpha w_j)}{T_{ij}}},
\]
где $q_{sj}$ --- часть времени сессии $s$, приписанная задаче $j$
(так что $T_{ij} = \sum_s q_{sj}$), а «решена в $s$» означает, что решающая
посылка попала в эту сессию.

\textbf{Читается так:} занятие месячной давности весит вдвое меньше
вчерашнего. Знаменатель --- та же семантика, что у базовой скорости: только
время решённых задач курса, кванты масштабированы тем же полом.

Показывается, только если знаменатель $\geq 60$ минут (эффективный час
недавнего продуктивного времени). Если ученик в последнее время лишь бился над
нерешённым, форма честно не показывается («---») вместо ложно низкого числа.

\section{Занятость и прогноз}

\textbf{$\approx$ч/нед.} Активное время на задачах курса раскладывается по
последним 8 неделям; берутся только недели, где активность была:
\[
  h_i \;=\; \med\bigl\{\, \text{часы недели} \;:\; \text{часы} > 0 \,\bigr\}.
\]
Нужно хотя бы 2 такие недели, иначе «---». Медиана, а не среднее: одна
пропущенная неделя не должна вдвое занижать оценку занятости.

\textbf{$\approx$недель.} Пусть $W_{\text{ост}} = \sum_{j \notin S_i} w_j$ ---
остаток курса по весу. Скорость считается по \emph{продуктивному} времени, а
будущие занятия опять будут содержать время на нерешаемое, поэтому остаток
масштабируется историческим КПД:
\[
  E_i \;=\; \min\Bigl(1,\
    \frac{\sum_{j \in S_i^{+}} \max(T_{ij},\,\alpha w_j)}{A_i}\Bigr)
  \qquad \text{(доля времени, ушедшая в решённые задачи).}
\]
Тогда
\[
  \text{осталось часов} \;\approx\; \frac{W_{\text{ост}}}{60\, v_i^{\text{тек}} E_i},
  \qquad
  \text{осталось недель} \;\approx\; \frac{W_{\text{ост}}}{60\, v_i^{\text{тек}} E_i\, h_i}.
\]
Без $E_i$ прогноз был бы систематически оптимистичным.

\textbf{Важно:} здесь используется \emph{сырая} $v_i^{\text{тек}}$ (до деления
на медиану группы) --- иначе размерность не сошлась бы. Прогноз считается
раньше нормировки.

\section{Сигналы}

Оба сигнала --- только про задачи, которые ученик \emph{пробовал, но не решил}.

\textbf{Застрял.}
\[
  z_{ij} = \frac{T_{ij}}{w_j} \;>\; z^{*} = 3.
\]
«Потратил уже втрое больше типичного времени и не решил» --- повод подойти.
Сортируются по убыванию $z_{ij}$.

\textbf{Пропустил.} Задача с попытками, после которой ученик решил
$\geq 2$ более поздних задачи курса: ушёл дальше, не закрыв. Сортируются по
порядку курса.

Каждого вида показывается не больше 4 штук.

\section{Флаги честности}
\label{sec:flags}

Четыре детектора подозрительных паттернов в посылках. Это \emph{не вердикт}, а
повод проверить лично.

\begin{center}
\begin{tabular}{@{}lp{9.6cm}@{}}
\toprule
Детектор & Срабатывает, когда \\
\midrule
Серия «с первой попытки» &
  5 подряд решённых с первой посылки \emph{нелёгких} задач (по когорте доля
  first-try ниже 0,6 при $\geq 5$ решивших); перерыв больше 14 дней рвёт серию \\
«Пулемёт» &
  $\geq 4$ решений подряд с паузами $\leq 3$ мин при суммарном типичном
  времени $\geq 40$ мин \\
«Резко быстрее» &
  окно из 5 подряд решённых, где личное время меньше 10\% типичного, при весе
  окна $\geq 30$ мин \\
«Пачечная сдача» &
  $\geq 4$ задачи курса решены в одной сессии при медианном личном времени
  меньше 50\% типичного (решения написаны заранее, сдана очередь посылок) \\
\bottomrule
\end{tabular}
\end{center}

\textbf{Как флаг влияет на темп.} По умолчанию эпизоду не доверяют: до
разметки преподавателем его посылки \emph{исключаются из всей математики}
темпа --- времён, скоростей и весов когорты. Решённые задачи при этом остаются
решёнными (прогресс не страдает), просто без времени --- как задачи ACMP.
Разметка: «сам решил» возвращает эпизод в подсчёт, «перенос» и «нарушение»
оставляют исключённым навсегда.

Считается это в две фазы: сначала на данных без эпизодов «перенос»/«нарушение»
детектируются сами флаги (их и видно в интерфейсе), затем из данных
дополнительно убираются неразмеченные эпизоды --- и по ним считается темп.

\section{Параметры}

\begin{center}
\begin{tabular}{@{}llp{7.4cm}@{}}
\toprule
Параметр & Значение & Смысл \\
\midrule
$\tau$ & 45 мин & порог разрыва сессии \\
$\delta_0$ & $\leq 10$ мин & надбавка на «вход» в сессию \\
$\alpha$ & $1/2$ & пол времени решённой задачи; ограничивает её вклад величиной $\times 2$ \\
$n_0$ & 5 & сила сглаживания весов к типичной задаче \\
$\bar w$ по умолчанию & 20 мин & если данных о времени нет вовсе \\
$\beta$, $c_{\max}$ & $1/2$, 3 & поправка весов на редкость решения \\
дошедших для поправки & $\geq 5$ & иначе доля слишком шумная \\
решивших для показа $\widehat w_j$ & $\geq 3$ & порог «типичного времени» в таблице задач \\
$H$ & 28 дней & полупериод забывания для текущей формы \\
порог формы & 60 мин & эффективное недавнее продуктивное время \\
$z^{*}$ & 3 & порог сигнала «застрял» \\
порог показа скорости & 5 задач, 2 ч & иначе «мало данных» \\
когорта для нормировки & $\geq 5$ учеников & иначе сырая шкала \\
недель для $h_i$ & $\geq 2$ активных & окно --- последние 8 недель \\
сигналов на ученика & $\leq 4$ & каждого вида \\
\bottomrule
\end{tabular}
\end{center}

\section{Границы применимости}

\begin{itemize}[itemsep=2pt]
  \item Учитываются только сайты со временем посылок (informatics, codeforces,
        ejudge). Задачи ACMP входят в прогресс, но не в скорость: если курс во
        многом на ACMP, скорость описывает лишь часть материала.
  \item «Активное время» --- \textbf{нижняя оценка}: обдумывание без посылок
        модели не видно. Поэтому метрика сравнительная (все ученики смещены
        одинаково), а не абсолютная: «в полтора раза быстрее группы» --- да,
        «решает за 20 минут» --- нет.
  \item Решение, написанное заранее и отправленное пачкой, завышает скорость.
        Пол $\alpha w_j$ ограничивает завышение по каждой задаче величиной
        $\times 2$, серии и пачки ловятся флагами честности и исключаются из
        темпа до разметки, остальное гасится медианами.
  \item Время, утопленное в нерешённых задачах, в скорость не входит ---
        осознанный выбор: скорость отвечает на вопрос «как быстро ученик
        проходит то, что проходит», а затыки показываются отдельно.
  \item В первые недели нового курса решивших мало, веса почти равны
        ($w_j \approx \bar w$), и скорость временно ближе к «задач в час».
  \item Прогноз ($\approx$недель) --- грубая оценка: он предполагает, что
        текущая форма и занятость сохранятся, а оставшиеся задачи по сложности
        похожи на пройденные.
\end{itemize}

\section{Где это в коде}

Всё описанное живёт в одном файле --- \texttt{course\_stats.go} в каталоге
\texttt{internal/standings/}:

\begin{center}
\begin{tabular}{@{}ll@{}}
\toprule
Что & Функция \\
\midrule
Сессии и $T_{ij}$ & \texttt{buildStudentTaskTimes} \\
Веса $w_j$ & \texttt{courseWeights} \\
«Типичное время» в таблице задач & \texttt{courseDisplayWeights} \\
Прогресс, скорости, прогноз, сигналы & \texttt{computeStudentCourseStats} \\
Нормировка на когорту & \texttt{normalizeCourseSpeeds} \\
Флаги честности & \texttt{detectCourseFlags} \\
Порядок задач курса & \texttt{courseTasksFromStandings} \\
\bottomrule
\end{tabular}
\end{center}

Параметры собраны константами в начале файла. Выбор когорты (группа или все
группы с этими контестами) делается рядом, в \texttt{builder.go} ---
функция \texttt{BuildGroupsStandings}.

\end{document}
